\chapter{An Opening Word}
Many voices had been speaking of eternal life, before the days of
the Son of man.  Especially pronounced had been the teachings of
the Egyptians that there is another world.   In their Acadian hymns
the Chaldaeans had dimly foretold a future life.  The belief of the
Parsees, as expressed in their Zend-Avesta, had included a place of
darkness for the evil soul and a reward for the good in the realm
of light.  The Hindus had declared, in their Rig-Veda, their
beautiful conception of the immortality of the soul, and had
written of a future "imperishable world, where there is eternal
light and glory."  The Grecian and Roman mythologies had voiced
their hope of blessedness for the shades of the departed.
Everywhere serious men had been asking as to the experiences beyond
the grave.  It was as if the Eastern world had become a vast
parliament chamber, wherein the nations were proclaiming their
different doctrines as to a future life.

In the midst of these varying and uncertain voices, Christ spoke
his authoritative message.  There was no wavering in his tone.
What the Oriental philosophers were guessing, he revealed; what the
Hebrew prophets had foreshadowed in their holy writings, he
unfolded in full light.  The ancient Vedic hymns, the oracles of
Greece, the Egyptian \textit{Book of the Dead}, anticipating by two
thousand years the Hebrew exodus--all these are naught compared
with the words of that inspired Teacher who spoke in Palestine.

In addition, Christ was himself the vital evidence of the
resurrection which he taught.  Against the assaults of doubt his
unique teachings are buttressed forevermore by his own return from
the land of silence.  In a short week after his words to Martha at
Bethany he had become, through his own rare experience, the
resurrection and the life.  Not the dead Buddha, nor the departed
Zoroaster, nor the vanished  Pythagoras  ever came back through the
opened door of the sepulcher, wearing the grave clothes of those
who sleep.  Human fancy had never dreamed of such a rapturous
denouement for faiths other than Christianity.  The resurrection of
the Lord is the crowning narrative with which the Gospels close.
It is a risen Christ who repairs the wastage of human decay and
death.  A voice above all those from Ind or Persia or the Nile
speaks henceforth in Judaea and the world concerning immortality.
The superlative Easter argument is the risen Christ himself.
